
\documentclass[11pt]{article}
\usepackage{amsmath,amssymb,amsthm,mathtools}
\usepackage[margin=1in]{geometry}

\newtheorem{theorem}{Theorem}
\newtheorem{lemma}{Lemma}
\newtheorem{proposition}{Proposition}
\newtheorem{conjecture}{Conjecture}
\newtheorem{definition}{Definition}

\title{A Rank-Two Affine S-Unit Prime-Production Route for Erdős--Graham 203}
\author{Jared Wilder}
\date{\today}

\begin{document}
\maketitle

\begin{abstract}
We isolate a finite analytic route to EG203.  The problem is reduced to a
rank-two affine S-unit prime-production theorem: for every \(m\) with
\((m,6)=1\), some \(m2^k3^\ell+1\) is prime.  The elementary Type-I and
Type-II algebraic reductions are recorded, and the remaining analytic input is
identified as a Vaughan/Heath-Brown decomposition, a rank-two S-unit
Type-II bound, a subgroup concentration estimate, and a Brun--Hooley
lower-bound sieve.
\end{abstract}

\section{Target theorem}

\begin{theorem}[EG203 prime production]
For every integer \(m\ge 1\) with \((m,6)=1\), there exist \(k,\ell\ge 0\)
such that
\[
  m2^k3^\ell+1
\]
is prime.
\end{theorem}

\section{Algebraic reductions}

\begin{lemma}[Type-I congruence]
If \(p\mid mA+1\), then
\[
  mA\equiv -1\pmod p.
\]
If \(p\) is prime, then \(p\nmid m\).
\end{lemma}

\begin{lemma}[Pair-GCD/S-unit difference]
If
\[
d\mid m2^k3^\ell+1,\qquad d\mid m2^{k'}3^{\ell'}+1,\qquad (d,m)=1,
\]
then
\[
d\mid 2^k3^\ell-2^{k'}3^{\ell'}.
\]
\end{lemma}

\begin{proof}
Subtract the two congruences and invert \(m\pmod d\).
\end{proof}

\section{Collision identity}

Let \(T_D=\{(k,\ell):k+\ell\le D\}\). Define
\[
C(D,q)=\#\{(x,y)\in T_D^2:2^{x_1}3^{x_2}\equiv2^{y_1}3^{y_2}\pmod q\}.
\]
For \(h=(a,b)\ne(0,0)\), let
\[
R_{a,b}=\left|2^{a_+}3^{b_+}-2^{a_-}3^{b_-}\right|.
\]
Let
\[
N_D(a,b)=\#\{(x,y)\in T_D^2:x-y=(a,b)\}.
\]
Then
\[
C(D,q)-|T_D|
=
\sum_{(a,b)\ne(0,0)}N_D(a,b)\mathbf 1_{q\mid R_{a,b}}.
\]

Moreover
\[
N_D(a,b)=
\begin{cases}
0,& S<0,\\
(S+1)(S+2)/2,& S\ge 0,
\end{cases}
\]
where
\[
S=\min(D,D-a-b)-\max(0,-a)-\max(0,-b).
\]

\section{Analytic theorem needed}

\begin{theorem}[Rank-two S-unit prime lower bound]
For every \(m\) with \((m,6)=1\), there exists \(D=D(m)\) such that
\[
\#\{(k,\ell)\in T_D:m2^k3^\ell+1\text{ prime}\}>0.
\]
\end{theorem}

\section{Required ingredients}

\subsection{Vaughan/Heath-Brown identity}
Apply a prime-detecting decomposition to
\[
\Lambda(m2^k3^\ell+1)
\]
over \(T_D\).

\subsection{Type I}
Use the congruence
\[
m2^k3^\ell\equiv -1\pmod p
\]
to control linear sums over the rank-two orbit.

\subsection{Type II}
By the pair-GCD lemma, off-diagonal common divisors reduce to
\[
2^k3^\ell-2^{k'}3^{\ell'}.
\]
The collision identity above reduces the problem to weighted divisor sums over
\(R_{a,b}\).

\subsection{Short-relation split}
Short relation vectors \((a,b)\) satisfying \(q\mid R_{a,b}\) must be handled
by their exact multiplicities \(N_D(a,b)\).

\subsection{Concentration estimate}
A usable form is
\[
\sum_{p\le z}\frac{1}{H_p}\ll \log\log z+O(1),
\]
where \(H_p=|\langle 2,3\rangle\subset\mathbb F_p^\times|\), with the precise
uniformity required by the sieve.

\subsection{Brun--Hooley lower-bound sieve}
Apply a lower-bound sieve at level \(r\ge \omega(z)+1\) to break parity and
obtain a positive count of prime values.

\section{Conclusion}
The elementary and combinatorial reductions are finite.  The remaining theorem
is the analytic rank-two S-unit prime lower bound above.
\end{document}
